% main.tex — INITIAL DRAFT written with AI assistance (Claude Code) at the
% author's request on 2026-09-12, to be rewritten by the author. Every number
% is a macro from numbers.tex (generated by the targets pipeline); tables and
% figures are generated files. See ai_disclosure.tex for the disclosure that
% this drafting assistance requires.
\documentclass[12pt,a4paper]{article}
\usepackage[utf8]{inputenc}
\usepackage[T1]{fontenc}
\usepackage{booktabs,graphicx,amsmath,natbib,hyperref,geometry,setspace}
\usepackage[flushleft]{threeparttable}
\geometry{margin=2.5cm}
\singlespacing
\input{numbers}

\title{Marine Heatwaves, Ocean Productivity and Fishing Effort:\\
High-Frequency Evidence from the Brazilian Coast}
\author{Valber Gregory\thanks{% AUTHOR WRITES affiliation, e-mail, acknowledgements.
}}
\date{Draft — \today}

\begin{document}
\maketitle

\begin{abstract}
Marine heatwaves (MHWs) are becoming longer and more frequent, yet evidence on how fishing fleets respond to them at high spatial and temporal resolution remains scarce, especially for tropical and subtropical fisheries. We assemble a reproducible geospatial data cube that links satellite-based apparent fishing effort (Global Fishing Watch) to daily sea-surface temperature (NOAA OISST v2.1) for the entire Brazilian coast and exclusive economic zone, at the level of \panelCells{} hexagonal cells (H3 resolution \gridRes{}) observed monthly from \panelStart{} to \panelEnd{} (\panelObs{} cell-months). MHWs follow the definition of Hobday et al.\ (2016). Poisson pseudo-maximum-likelihood models with cell, calendar-month and region-by-season fixed effects and Conley standard errors indicate that each additional MHW day in a month is associated with a change of \bMhwMain{} log points in fishing hours (about $-1.1\%$), conditional on the sea-surface temperature anomaly, which itself attracts effort (\bSstAnom{} per~$^{\circ}$C). The response operates on the intensive margin, is absent before exposure, is concentrated in months with twenty or more MHW days (about \pctBinTwentyPlus{}\%), and is robust to grid resolution, alternative MHW definitions, Conley cut-offs and sample exclusions. A specification with cell-specific linear trends, which also neutralises a twelve-month-lead placebo, yields a more conservative \bMhwCellTrend{}. We find no short-distance displacement but a positive association with neighbours' exposure at 60--150~km, and the largest responses for drifting longliners, offshore cells and the North region. The data cube is released as an open information-systems artefact for climate-resilient fisheries monitoring.
\end{abstract}

\noindent\textbf{Keywords:} marine heatwaves; fishing effort; AIS; spatial panel; Brazil; data cube\\
\noindent\textbf{JEL:} Q22; Q54; C23; R12

\section{Introduction}\label{sec:intro}
Marine heatwaves --- discrete periods of anomalously warm sea-surface temperature (SST) lasting at least five days \citep{hobday2016} --- have become longer and more frequent over the past century \citep{oliver2018} and now threaten marine biodiversity and the ecosystem services on which fisheries depend \citep{smale2019}. Historical warming has already reduced the productivity of many exploited stocks \citep{free2019}, single extreme events can reorganise fishing seasons and markets, as the 2012 Northwest Atlantic heatwave showed \citep{mills2013}, and projections indicate that heatwaves compound the losses caused by gradual warming \citep{cheung2020}. Understanding how fleets respond while a heatwave unfolds is therefore a first-order question for fisheries economics and for climate adaptation policy.

Most of the available evidence concerns temperate stocks and relies on landings or stock assessments at annual or seasonal resolution. Climate shocks that trigger closures impose costs that are mediated by the diversity of fishing portfolios \citep{fisher2021}, and fleets are known to reallocate effort laterally when part of their habitat becomes unusable, the canonical case being the Gulf of Mexico hypoxic zone \citep{smith2014}. Recent work with vessel-tracking data shows that MHWs redistribute pelagic fleets at basin scale, and that the spatial extent of a heatwave matters more than its intensity for the contraction of fishing grounds \citep{farchadi2024}. Much less is known, however, about the high-frequency, cell-level response of a tropical fleet, and almost nothing about the South Atlantic, where heatwaves have a documented regional driver --- the atmospheric blocking that simultaneously causes South American droughts \citep{rodrigues2019}. Brazil is a demanding case for a further reason: it has had no continuous national fisheries statistics since the official series were discontinued in the 2010s, so behavioural responses cannot be studied from landings at all. % AUTHOR: add one or two sentences on the institutional context.

This paper exploits the fact that satellite tracking of vessels through the Automatic Identification System (AIS) now makes fishing activity observable at fine spatial and temporal scales \citep{kroodsma2018}. We combine twelve years of Global Fishing Watch apparent fishing effort with daily NOAA OISST temperatures for the whole Brazilian coast in a cell-by-month panel and ask how MHW exposure is associated with the intensity and the spatial distribution of effort. The empirical strategy exploits within-cell variation in exposure, net of calendar-month shocks and regional seasonality, using Poisson pseudo-maximum-likelihood estimation \citep{santossilva2006,correia2020} and spatially robust inference \citep{conley1999}. Because AIS coverage is skewed towards larger vessels and apparent fishing hours are a classifier output \citep{hintzen2025}, we state explicitly that all inferences concern the AIS-tracked fleet, which in Brazil is dominated by drifting longliners.

Three findings emerge. First, MHW exposure is associated with fewer fishing hours in the affected cell, but only once the ordinary SST anomaly is held constant: a warm-but-not-extreme month attracts effort, whereas each MHW day is associated with \bMhwMain{} log points of hours, a response that operates almost entirely on the intensive margin and is concentrated in months with twenty or more heatwave days. Second, the response is dynamic and shows no anticipation: leads at one to three months are null, the contemporaneous and second-lag coefficients are negative, and the cumulative effect over six months is \bCumSixIC{} log points per MHW day. Third, the response is heterogeneous: it is largest in the North region, offshore and for drifting longliners, and there is no displacement to the immediately adjacent ring but a positive association with neighbours' exposure at 60--150~km. We treat the last result, and the causal reading of all of them, with caution: a twelve-month-lead placebo reveals differential trends across cells, and the specification that neutralises it halves the main coefficient to \bMhwCellTrend{}.

The paper makes three contributions. It provides the first coast-wide, high-frequency evidence on fleet responses to marine heatwaves for Brazil and, to our knowledge, for a tropical South Atlantic fishery. It offers a transparent treatment of the identification threats --- anticipation, differential trends and spatial dependence --- with placebos that the reader can inspect. And it delivers the Marine Fisheries Data Cube, a reproducible pipeline and storage design that integrates effort, temperature, heatwave events and territorial layers, as an information-systems contribution for monitoring and adaptation in a country without official fisheries statistics.

\section{Research context}\label{sec:context}
% AUTHOR WRITES: fragmentation of Brazilian fisheries data; discontinued national statistics; management institutions (MPA/IBAMA, closed seasons, Seguro-Defeso); the tuna and swordfish longline fleet; AIS mandates.

\section{Data}\label{sec:data}
\paragraph{Fishing effort.} Global Fishing Watch (GFW) infers fishing behaviour from AIS tracks with a neural-network classifier and publishes ``apparent fishing hours'' through its 4Wings API, aggregated monthly on a $0.1^{\circ}$ grid by flag and gear type \citep{kroodsma2018}. We query the box $54^{\circ}$W--$25^{\circ}$W, $35^{\circ}$S--$6^{\circ}$N for \panelStart{}--\panelEnd{} (twelve complete years; 2012, the ramp-up year of AIS coverage, is excluded) and map each $0.1^{\circ}$ grid centre to an H3 hexagonal cell. Two limitations must be stated at the outset. AIS coverage is concentrated in larger vessels; the artisanal fleet, which accounts for a large share of Brazilian landings, is essentially absent, so the population under study is the AIS-tracked fleet. Apparent fishing hours are also a model output: validation against on-board records for Northeast Atlantic pelagic fleets shows that slow-speed searching can be classified as fishing, overstating gear-deployment time by fleet-specific factors \citep{hintzen2025}. Cell fixed effects absorb the time-invariant part of this error, and gear-specific estimates address its fleet-dependent part, but the outcome should be read as inferred fishing activity rather than measured gear time.

\paragraph{Sea-surface temperature and marine heatwaves.} NOAA OISST v2.1 provides daily SST on a $0.25^{\circ}$ grid from September 1981. For each ocean pixel we compute a \climStart{}--\climEnd{} daily climatology and the seasonally varying \mhwPctile{}th-percentile threshold (31-day smoothing window), and detect MHWs as runs of at least \mhwMinDuration{} consecutive days above the threshold, merging gaps of at most two days \citep{hobday2016,hobday2018}; the implementation is the \texttt{heatwaveR} package \citep{schlegel2018}. The monthly exposure variables are the number of MHW days in the month (\texttt{mhw\_days}), the maximum intensity above the threshold, the mean SST and the mean anomaly relative to the seasonal climatology (\texttt{sst\_anom}). The anomaly is kept as a separate regressor because a warm anomaly is not a heatwave; the two enter jointly in every model.

\paragraph{Grid, universe and territorial layers.} Cells are H3 hexagons of resolution \gridRes{} (about 1,770~km$^2$), chosen after comparing resolutions 4 and 5 on the share of zeros and the within-cell variance of effort; resolution 5 is retained as a robustness check. The analysis universe is the fishing footprint, that is, the cells with positive effort in at least one month of the sample, giving \panelCells{} cells $\times$ 144 months $=$ \panelObs{} cell-months, of which \shareEffortPos{}\% have positive hours and \shareMhwAny{}\% have at least one MHW day (mean \meanMhwDays{} days per cell-month). Each cell carries its distance to the coast and the nearest coastal macro-region (North, Northeast, Southeast, South) from IBGE boundaries. Table~\ref{tab:desc} reports descriptive statistics by region; Figure~\ref{fig:maps} maps mean effort and mean MHW exposure; Figure~\ref{fig:ts} shows the national time series of hours and of the share of cells under a heatwave.

\begin{table}[htbp]
\centering
\caption{Descriptive statistics by coastal region, 2013-2024 (H3 res.\ 4 fishing footprint)}
\label{tab:desc}
\begin{tabular}{lrrrrrrrrr}
\hline
region & cells & cell\_months & share\_pos & hours\_total & hours\_mean\_pos & sst\_mean & mhw\_days\_mean & share\_mhw\_any & dist\_coast\_km
\\ \hline
Sul &  650 &  93600 & 0.330 & 4090253 & 132.4 & 21.21 & 4.28 & 0.344 &  439 \\
Sudeste & 1916 & 275904 & 0.159 & 2807945 &  64.0 & 23.84 & 4.16 & 0.295 & 1019 \\
Nordeste & 2001 & 288144 & 0.119 &  809156 &  23.6 & 27.61 & 5.59 & 0.367 &  618 \\
Norte &  198 &  28512 & 0.108 &  447580 & 145.6 & 28.22 & 5.79 & 0.420 &  268 \\
All & 4765 & 686160 & 0.163 & 8154935 &  72.7 & 25.25 & 4.84 & 0.337 &  717 \\
\hline
\end{tabular}
\end{table}


\begin{figure}[htbp]\centering
\includegraphics[width=\linewidth]{figures/fig_maps_effort_mhw.png}
\caption{Mean monthly apparent fishing hours by cell (left, log scale) and mean marine-heatwave days per month (right), Brazilian coast and EEZ, \panelStart{}--\panelEnd{}.}\label{fig:maps}
\end{figure}

\begin{figure}[htbp]\centering
\includegraphics[width=0.9\linewidth]{figures/fig_timeseries.png}
\caption{Total monthly apparent fishing hours (bars) and share of cells under a marine heatwave (line), national aggregate.}\label{fig:ts}
\end{figure}

\section{The Marine Fisheries Data Cube}\label{sec:cube}
% AUTHOR WRITES (IS contribution). Draft outline: the cube is a targets pipeline (R) with 60+ targets that download raw files with checksums, derive climatologies and MHW events per pixel with block-level checkpoints, build H3 grids, assemble cell-month panels stored as Parquet and queried with DuckDB, run the estimation and export every table, figure and number as generated files; lineage is the pipeline graph; people--process--technology framing.

\section{Empirical strategy}\label{sec:methods}
\paragraph{Outcome and estimator.} Let $y_{it}$ denote apparent fishing hours in cell $i$ and month $t$. Because \shareEffortPos{}\% of cell-months are positive and the remainder are exact zeros, a log-linear model is unsuitable: $\log(1+y)$ is scale-dependent and inconsistent under heteroskedasticity \citep{santossilva2006}. We therefore model the conditional mean by Poisson pseudo-maximum likelihood,
\begin{equation}\label{eq:main}
E[y_{it}\mid X_{it}] = \exp\!\big(\beta\,\mathrm{MHW}_{it} + \gamma\,\mathrm{ANOM}_{it} + \alpha_i + \delta_t + \theta_{r(i),m(t)}\big),
\end{equation}
where $\mathrm{MHW}_{it}$ is the number of MHW days, $\mathrm{ANOM}_{it}$ the mean SST anomaly, $\alpha_i$ a cell fixed effect, $\delta_t$ a calendar-month fixed effect (144 levels) and $\theta_{r(i),m(t)}$ a region $\times$ month-of-year effect. PPML is consistent whenever the conditional mean is correctly specified, regardless of the distribution of $y$, and $100\beta$ is read as the percentage change in expected hours per additional MHW day. Estimation uses the iteratively reweighted algorithm for high-dimensional fixed effects \citep{correia2020}, implemented in the R package \texttt{fixest}.

\paragraph{What the fixed effects absorb.} The cell effect $\alpha_i$ removes every time-invariant determinant of effort --- depth, distance to ports, average productivity, permanent closures, the cell's average AIS coverage. The calendar-month effect $\delta_t$ removes shocks common to all cells in a month --- fuel prices, demand, the average state of ENSO, national changes in AIS adoption. The region $\times$ month-of-year effect removes region-specific seasonality such as annual closed seasons. Identification of $\beta$ therefore rests on within-cell deviations of MHW exposure from the cell's own mean, net of the month's national shock and the region's seasonal norm. We use associational language throughout: the design does not randomise exposure, and Section~\ref{sec:rob} shows that differential trends across cells are a real threat.

\paragraph{Inference.} Neighbouring cells share the same SST pixel and the same vessels, and months are serially dependent. We report \citet{conley1999} standard errors with a \conleyCutoff{}~km spatial cut-off (100 and 400~km as robustness), which also accommodate serial correlation within cell; a two-way cluster by cell and month is the fallback. Moran's $I$ of cell-mean residuals is \moranI{} (permutation $p=\pMoran{}$), confirming small but non-zero residual spatial dependence and justifying this choice.

\paragraph{Margins and dynamics.} The extensive margin uses $\mathbf{1}[y_{it}>0]$ as the outcome in a linear probability model with the same fixed effects. Dynamics use distributed leads and lags, $\sum_{k=1}^{3}\phi_k\,\mathrm{MHW}_{i,t+k} + \sum_{k=0}^{6}\beta_k\,\mathrm{MHW}_{i,t-k}$; under no anticipation the lead coefficients $\phi_k$ should be zero, and the cumulative effect $\sum_k\beta_k$ is reported with its Conley variance. Non-linearity replaces the linear index by indicators for bins of MHW days (0 as reference, 1--4, 5--9, 10--19, 20 or more) and of the anomaly.

\paragraph{Spatial displacement.} For each cell we compute the mean MHW exposure of its neighbours in three concentric rings defined by centroid distance --- 0--60~km (the first hexagonal ring), 60--150~km and 150--300~km --- using a row-normalised sparse contiguity matrix, and add them to equation~\eqref{eq:main}. A positive ring coefficient, conditional on own exposure, indicates that effort moves into the cell when its surroundings are under a heatwave, the mechanism documented for hypoxia by \citet{smith2014}. At the region $\times$ month level we also examine the effort-weighted mean distance to coast, a Herfindahl index of the spatial concentration of hours and the number of active cells, using historically fixed cell weights for the regional exposure measure so that the regressor is not contaminated by the contemporaneous composition of effort.

\paragraph{Heterogeneity.} Interactions of MHW days with region, with distance to coast (above and below the median) and, in a cell $\times$ month $\times$ gear panel with cell $\times$ gear fixed effects, with gear type.

\paragraph{Threats and placebos.} (i) Anticipation: leads at one to three months. (ii) Differential trends: MHW exposure rises steeply over the period, from below one to above thirteen days per cell-month between 2013 and 2024, so cells with growing exposure may also have declining effort for unrelated reasons; we add region $\times$ year effects and, more demandingly, cell-specific linear trends, and we use a twelve-month lead conditional on lags 0--6 as the diagnostic. (iii) Spatial confounding: we permute the MHW series across cells within region \nPermutations{} times and compare the observed coefficient with the permutation distribution. (iv) Measurement of MHWs: 95th percentile, ten-day minimum duration and a 1982--2011 climatology. (v) Sample: the resolution-5 grid, cells active in at least twelve months, and exclusion of cells within 25~km of the coast (a proxy for port areas) and of low-coverage cells.

\section{Results}\label{sec:results}
\subsection{Main estimates}
Table~\ref{tab:main} reports the baseline. Without the anomaly control, the MHW coefficient is small and imprecise (\bMhwNoControl{}, $p=\pMhwNoControl{}$). Conditioning on the anomaly separates two opposite associations that the unconditional estimate confounds: a warm-but-not-extreme month attracts effort (\bSstAnom{} log points per~$^{\circ}$C, SE \seSstAnom{}), whereas each MHW day is associated with \bMhwMain{} log points (SE \seMhwMain{}, $p=\pMhwMain{}$), about $-1.1\%$ of expected hours. The two coefficients together describe a non-monotonic thermal response: moderate warming improves catchability or accessibility for a pelagic fleet that targets warm-water species, while sustained extremes reduce activity. The extensive margin does not respond (\bExtensive{}, $p=\pExtensive{}$): vessels do not abandon cells, they fish fewer hours in them. The $\log(1+y)$ specification gives the same sign with a smaller magnitude (\bLogOnePlus{}), as expected from the compression of that transformation, and the maximum-intensity term is positive but imprecise once days are controlled (\bMaxIntensity{}, $p=\pMaxIntensity{}$), indicating that duration rather than peak intensity carries the association.


\begin{table}[htbp]
   \caption{\label{tab:main} Marine heatwave exposure and apparent fishing effort}
   \centering
   \begin{tabular}{lcccccc}
      \tabularnewline \midrule \midrule
      Dependent Variables: & \multicolumn{2}{c}{Fishing hours} & Any effort (0/1) & log1p(hours) & \multicolumn{2}{c}{Fishing hours}\\
      Model:               & (1)          & (2)             & (3)       & (4)           & (5)             & (6)\\  
                           &  Poisson     & Poisson         & OLS       & OLS           & Poisson         & Poisson\\  
      \midrule
      \emph{Variables}\\
      MHW days             & -0.0041      & -0.0113$^{***}$ & -0.0006   & -0.0033$^{*}$ &                 & -0.0068$^{*}$\\   
                           & (0.0034)     & (0.0035)        & (0.0004)  & (0.0018)      &                 & (0.0037)\\   
      SST anomaly          &              & 0.1308$^{***}$  &           &               &                 &   \\   
                           &              & (0.0398)        &           &               &                 &   \\   
      l(mhw\_days, 0)      &              &                 &           &               & -0.0018         &   \\   
                           &              &                 &           &               & (0.0020)        &   \\   
      l(mhw\_days, 1)      &              &                 &           &               & -0.0009         &   \\   
                           &              &                 &           &               & (0.0025)        &   \\   
      l(mhw\_days, 2)      &              &                 &           &               & -0.0059$^{***}$ &   \\   
                           &              &                 &           &               & (0.0019)        &   \\   
      l(mhw\_days, 3)      &              &                 &           &               & -0.0003         &   \\   
                           &              &                 &           &               & (0.0026)        &   \\   
      MHW max. intensity   &              &                 &           &               &                 & 0.0607$^{*}$\\   
                           &              &                 &           &               &                 & (0.0314)\\   
      \midrule
      \emph{Fixed-effects}\\
      Cell                 & Yes          & Yes             & Yes       & Yes           & Yes             & Yes\\  
      Month                & Yes          & Yes             & Yes       & Yes           & Yes             & Yes\\  
      Region-Month of year & Yes          & Yes             & Yes       & Yes           & Yes             & Yes\\  
      \midrule
      \emph{Fit statistics}\\
      Observations         & 686,160      & 686,160         & 686,160   & 686,160       & 671,724         & 686,160\\  
      Squared Correlation  & 0.50441      & 0.50691         & 0.21383   & 0.32667       & 0.50530         & 0.50478\\  
      Pseudo R$^2$         & 0.73412      & 0.73454         & 0.28358   & 0.12324       & 0.73481         & 0.73415\\  
      BIC                  & 15,545,556.4 & 15,520,953.8    & 483,640.4 & 1,997,387.1   & 15,397,831.3    & 15,543,585.5\\  
      \midrule \midrule
      \multicolumn{7}{l}{\emph{Conley (200km) standard-errors in parentheses}}\\
      \multicolumn{7}{l}{\emph{Signif. Codes: ***: 0.01, **: 0.05, *: 0.1}}\\
   \end{tabular}
   
   \par \raggedright 
   Conley standard errors (200 km). Associational estimates; see identification strategy.
\end{table}




\subsection{Dynamics and anticipation}
Table~\ref{tab:dyn} and Figure~\ref{fig:event} show the distributed-lead-and-lag specification. The three leads are individually and jointly indistinguishable from zero (\bLeadThree{}, \bLeadTwo{} and \bLeadOne{} at $-3$, $-2$ and $-1$ months; largest $|t|=\maxAbsTLead{}$), so effort does not fall before the heatwave exists --- the standard no-anticipation placebo passes. The contemporaneous coefficient is \bLagZero{} ($p=\pLagZero{}$), the second lag \bLagTwo{} ($p=\pLagTwo{}$) and the fourth \bLagFour{} ($p=\pLagFour{}$); the response dissipates by the sixth month (\bLagSix{}, $p=\pLagSix{}$). The cumulative 0--6 month effect is \bCumSixIC{} log points per MHW day (95\% CI \loCumSix{} to \hiCumSix{}), so a ten-day heatwave is associated with roughly a fifth less effort distributed over the following half-year rather than a one-month dip. The alternation of the lag profile (significant at 0, 2 and 4 months, not at 1 and 3) is consistent with trip-length periodicity in the longline fleet and should not be over-interpreted.


\begin{table}[htbp]
   \caption{\label{tab:dyn} Dynamic specification: leads (placebo) and lags}
   \centering
   \begin{tabular}{lcc}
      \tabularnewline \midrule \midrule
      Dependent Variable: & \multicolumn{2}{c}{Fishing hours}\\
      Model:               & (1)                   & (2)\\  
      \midrule
      \emph{Variables}\\
      f(mhw\_days, 3)      & -0.0013               &   \\   
                           & (0.0023)              &   \\   
      f(mhw\_days, 2)      & -0.0014               &   \\   
                           & (0.0026)              &   \\   
      f(mhw\_days, 1)      & $4.42\times 10^{-6}$  &   \\   
                           & (0.0021)              &   \\   
      l(mhw\_days, 0)      & -0.0066$^{***}$       & -0.0085$^{***}$\\   
                           & (0.0023)              & (0.0024)\\   
      l(mhw\_days, 1)      & -0.0021               & -0.0014\\   
                           & (0.0025)              & (0.0024)\\   
      l(mhw\_days, 2)      & -0.0065$^{***}$       & -0.0053$^{***}$\\   
                           & (0.0018)              & (0.0018)\\   
      l(mhw\_days, 3)      & 0.0018                & 0.0031\\   
                           & (0.0025)              & (0.0024)\\   
      l(mhw\_days, 4)      & -0.0061$^{**}$        & -0.0046$^{**}$\\   
                           & (0.0024)              & (0.0023)\\   
      l(mhw\_days, 5)      & -0.0035$^{*}$         & -0.0025\\   
                           & (0.0020)              & (0.0018)\\   
      l(mhw\_days, 6)      & -0.0019               & -0.0030\\   
                           & (0.0018)              & (0.0019)\\   
      SST anomaly          & 0.1290$^{***}$        & 0.1324$^{***}$\\   
                           & (0.0431)              & (0.0374)\\   
      \midrule
      \emph{Fixed-effects}\\
      Cell                 & Yes                   & Yes\\  
      Month                & Yes                   & Yes\\  
      Region-Month of year & Yes                   & Yes\\  
      \midrule
      \emph{Fit statistics}\\
      Observations         & 641,790               & 657,432\\  
      Squared Correlation  & 0.50849               & 0.51027\\  
      Pseudo R$^2$         & 0.73486               & 0.73620\\  
      BIC                  & 14,581,096.5          & 15,208,609.3\\  
      \midrule \midrule
      \multicolumn{3}{l}{\emph{Conley (200km) standard-errors in parentheses}}\\
      \multicolumn{3}{l}{\emph{Signif. Codes: ***: 0.01, **: 0.05, *: 0.1}}\\
   \end{tabular}
\end{table}




\begin{figure}[htbp]\centering
\includegraphics[width=0.85\linewidth]{figures/fig_event_study_res4.png}
\caption{PPML coefficients on MHW days by horizon, $-3$ to $+6$ months, with 95\% Conley confidence intervals. Horizons left of the dotted line are leads (placebo).}\label{fig:event}
\end{figure}

\subsection{Non-linearity}
Table~\ref{tab:nonlin} and Figure~\ref{fig:dose} replace the linear index by bins. The response is a threshold rather than a gradient: months with 20 or more MHW days show \bBinTwentyPlus{} log points (about \pctBinTwentyPlus{}\%, $p=\pBinTwentyPlus{}$), months with 10--19 days \bBinTenNineteen{} ($p=\pBinTenNineteen{}$), and shorter exposures are indistinguishable from zero (\bBinFiveNine{} and \bBinOneFour{}). The anomaly bins are individually imprecise but positive between $+0.5$ and $+1^{\circ}$C, mirroring the positive linear anomaly coefficient. The linear coefficient of Table~\ref{tab:main} is therefore an average over a response that is essentially zero for most exposures and large when the month is almost entirely under a heatwave --- the situation that the trend in exposure makes increasingly common.

\begin{table}[htbp]
\centering
\caption{Non-linear response: bins of SST anomaly and of MHW days}
\label{tab:nonlin}
\begin{tabular}{lrrrr}
\hline
model & term & b & se & p
\\ \hline
bins_anom & anom_bin::lt_m1 & -0.1641 & 0.1528 & 0.2831 \\
bins_anom & anom_bin::m1_m05 &  0.0257 & 0.1002 & 0.7978 \\
bins_anom & anom_bin::0_05 &  0.0525 & 0.0525 & 0.3171 \\
bins_anom & anom_bin::05_1 &  0.1356 & 0.0999 & 0.1745 \\
bins_anom & anom_bin::1_15 &  0.0695 & 0.0953 & 0.4657 \\
bins_anom & anom_bin::gt_15 & -0.0336 & 0.1455 & 0.8172 \\
bins_mhw & mhw_bin::d1_4 & -0.0011 & 0.0713 & 0.9881 \\
bins_mhw & mhw_bin::d5_9 & -0.0131 & 0.0286 & 0.6457 \\
bins_mhw & mhw_bin::d10_19 & -0.0594 & 0.0455 & 0.1924 \\
bins_mhw & mhw_bin::d20p & -0.3281 & 0.0911 & 0.0003 \\
bins_mhw & sst_anom &  0.1098 & 0.0371 & 0.0030 \\
\hline
\end{tabular}
\end{table}


\begin{figure}[htbp]\centering
\includegraphics[width=\linewidth]{figures/fig_dose_response.png}
\caption{Coefficients for bins of SST anomaly and of MHW days relative to the reference bins, with 95\% Conley intervals.}\label{fig:dose}
\end{figure}

\section{Spatial displacement}\label{sec:spatial}
Table~\ref{tab:spill} adds the exposure of neighbours in three rings. Own exposure remains negative. The 0--60~km ring is negative and insignificant (\bRingNear{}, $p=\pRingNear{}$): a heatwave that hits a cell typically also hits its first ring, so there is no nearby refuge to move to. The 60--150~km ring is positive (\bRingMid{}, $p=\pRingMid{}$), consistent with relocation over one or two days of steaming, and the 150--300~km ring is null (\bRingFar{}, $p=\pRingFar{}$). Table~\ref{tab:reloc} examines the aggregate consequences at the region $\times$ month level: regional MHW exposure is not associated with the effort-weighted distance to coast (\bRelocDistCoast{}) or with spatial concentration (\bRelocHHI{}), but the number of active cells increases (\bRelocActive{}), so effort spreads over more cells when a region is exposed. These aggregate estimates rest on only \nRegions{} regions and should be read as descriptive. Taken together, the results support a qualified version of the displacement hypothesis of \citet{smith2014} and \citet{farchadi2024}: effort does not reallocate to adjacent cells, because adjacency does not provide thermal refuge in an event of regional extent, but it does spread to intermediate distances.


\begin{table}[htbp]
   \caption{\label{tab:spill} Spatial displacement: neighbourhood exposure by distance ring}
   \centering
   \begin{tabular}{lcc}
      \tabularnewline \midrule \midrule
      Dependent Variables:       & Fishing hours  & Any effort (0/1)\\  
      Model:                     & (1)            & (2)\\  
                                 &  Poisson       & OLS\\  
      \midrule
      \emph{Variables}\\
      MHW days                   & -0.0026        & 0.0002\\   
                                 & (0.0031)       & (0.0002)\\   
      Neighbour MHW (0-60 km)    & -0.0147        & -0.0002\\   
                                 & (0.0096)       & (0.0003)\\   
      Neighbour MHW (60-150 km)  & 0.0147$^{*}$   & 0.0010\\   
                                 & (0.0075)       & (0.0006)\\   
      Neighbour MHW (150-300 km) & -0.0122        & -0.0026$^{***}$\\   
                                 & (0.0098)       & (0.0009)\\   
      SST anomaly                & 0.1365$^{***}$ & 0.0044\\   
                                 & (0.0401)       & (0.0059)\\   
      \midrule
      \emph{Fixed-effects}\\
      Cell                       & Yes            & Yes\\  
      Month                      & Yes            & Yes\\  
      Region-Month of year       & Yes            & Yes\\  
      \midrule
      \emph{Fit statistics}\\
      Observations               & 686,160        & 686,160\\  
      Squared Correlation        & 0.50695        & 0.21417\\  
      Pseudo R$^2$               & 0.73468        & 0.28410\\  
      BIC                        & 15,512,704.7   & 483,391.9\\  
      \midrule \midrule
      \multicolumn{3}{l}{\emph{Conley (200km) standard-errors in parentheses}}\\
      \multicolumn{3}{l}{\emph{Signif. Codes: ***: 0.01, **: 0.05, *: 0.1}}\\
   \end{tabular}
\end{table}



\begin{table}[htbp]
\centering
\caption{Reallocation outcomes at the region-month level}
\label{tab:reloc}
\begin{tabular}{lrrrr}
\hline
model & term & b & se & p
\\ \hline
dist\_coast & mhw\_exp &  0.4428 &  2.4997 & 0.8595 \\
dist\_coast & sst\_anom &  7.7878 & 40.8349 & 0.8488 \\
concentration & mhw\_exp & -0.0012 &  0.0008 & 0.1260 \\
concentration & sst\_anom & -0.0017 &  0.0152 & 0.9115 \\
active\_cells & mhw\_exp &  0.0192 &  0.0045 & 0.0000 \\
active\_cells & sst\_anom & -0.2480 &  0.0826 & 0.0027 \\
\hline
\end{tabular}
\end{table}


\section{Heterogeneity}\label{sec:het}
Table~\ref{tab:het} interacts exposure with region, distance and gear. By region, the association is largest in the North (\bNorth{}, $p=\pNorth{}$), intermediate in the Southeast (\bSoutheast{}, $p=\pSoutheast{}$) and Northeast (\bNortheast{}, $p=\pNortheast{}$), and absent in the South (\bSouth{}, $p=\pSouth{}$) --- a tropical-to-subtropical gradient consistent with tropical species operating closer to their thermal limits. By distance, offshore cells respond (\bOffshore{}, $p=\pOffshore{}$) and nearshore cells do not (\bNearshore{}, $p=\pNearshore{}$). By gear, drifting longliners respond (\bDriftLongline{}, $p=\pDriftLongline{}$), whereas set longlines (\bSetLongline{}, $p=\pSetLongline{}$) and trawlers (\bTrawl{}, $p=\pTrawl{}$) are not distinguishable from zero. The response is thus concentrated in the offshore pelagic longline fleet in tropical waters, the segment for which a surface-temperature signal is most directly linked to the distribution of target species. % AUTHOR: add the fisheries-economics interpretation (tuna/swordfish thermal habitat, trip economics of the longline fleet).


\begin{table}[htbp]
   \caption{\label{tab:het} Heterogeneity by region, gear type and distance to coast}
   \centering
   \begin{tabular}{lccc}
      \tabularnewline \midrule \midrule
      Dependent Variable: & \multicolumn{3}{c}{Fishing hours}\\
      Model:                                           & (1)             & (2)             & (3)\\  
      \midrule
      \emph{Variables}\\
      MHW days $\times$ Region $=$ Nordeste            & -0.0120$^{**}$  &                 &   \\   
                                                       & (0.0056)        &                 &   \\   
      MHW days $\times$ Region $=$ Norte               & -0.0409$^{***}$ &                 &   \\   
                                                       & (0.0054)        &                 &   \\   
      MHW days $\times$ Region $=$ Sudeste             & -0.0108$^{***}$ &                 &   \\   
                                                       & (0.0038)        &                 &   \\   
      MHW days $\times$ Region $=$ Sul                 & 0.0041          &                 &   \\   
                                                       & (0.0047)        &                 &   \\   
      SST anomaly                                      & 0.0898$^{**}$   & 0.1250$^{***}$  & 0.1168$^{***}$\\   
                                                       & (0.0396)        & (0.0410)        & (0.0406)\\   
      MHW days $\times$ gear $=$ drifting\_longlines   &                 & -0.0126$^{***}$ &   \\   
                                                       &                 & (0.0037)        &   \\   
      MHW days $\times$ gear $=$ fishing               &                 & -0.0114$^{***}$ &   \\   
                                                       &                 & (0.0042)        &   \\   
      MHW days $\times$ gear $=$ other                 &                 & -0.0107$^{**}$  &   \\   
                                                       &                 & (0.0042)        &   \\   
      MHW days $\times$ gear $=$ set\_longlines        &                 & -0.0072         &   \\   
                                                       &                 & (0.0076)        &   \\   
      MHW days $\times$ gear $=$ trawlers              &                 & -0.0102         &   \\   
                                                       &                 & (0.0105)        &   \\   
      MHW days $\times$ far\_coast $=$ nearshore       &                 &                 & -0.0074\\   
                                                       &                 &                 & (0.0047)\\   
      MHW days $\times$ far\_coast $=$ offshore        &                 &                 & -0.0239$^{***}$\\   
                                                       &                 &                 & (0.0042)\\   
      \midrule
      \emph{Fixed-effects}\\
      Cell                                             & Yes             &                 & Yes\\  
      Month                                            & Yes             & Yes             & Yes\\  
      Region-Month of year                             & Yes             & Yes             & Yes\\  
      Cell x gear                                      &                 & Yes             & \\  
      \midrule
      \emph{Fit statistics}\\
      Observations                                     & 686,160         & 1,410,480       & 686,160\\  
      Squared Correlation                              & 0.51159         & 0.39574         & 0.50720\\  
      Pseudo R$^2$                                     & 0.73550         & 0.67756         & 0.73480\\  
      BIC                                              & 15,465,150.8    & 20,498,480.0    & 15,505,862.5\\  
      \midrule \midrule
      \multicolumn{4}{l}{\emph{Conley (200km) standard-errors in parentheses}}\\
      \multicolumn{4}{l}{\emph{Signif. Codes: ***: 0.01, **: 0.05, *: 0.1}}\\
   \end{tabular}
\end{table}




\section{Robustness and placebos}\label{sec:rob}
Table~\ref{tab:rob} shows that the baseline coefficient is stable across the resolution-5 grid (\bResFive{}), region $\times$ year trends (\bTrendRegYear{}), cells active in at least twelve months (\bStableCells{}), Conley cut-offs of 100 and 400~km (identical point estimates), and the three alternative MHW definitions: the 95th-percentile threshold gives \bMhwPninetyfive{}, the ten-day minimum duration \bMhwDurTen{} and the 1982--2011 climatology \bMhwClimEarly{}. The ordering --- larger effect for the more extreme definition --- is a dose-response pattern that a spurious association would not be expected to produce. Table~\ref{tab:excl} shows that excluding cells within 25~km of the coast or low-coverage cells leaves the estimate unchanged.


\begin{table}[htbp]
   \caption{\label{tab:rob} Robustness: grid resolution, trends, sample, Conley cutoffs and MHW definition}
   \centering
   \begin{tabular}{lcccccc}
      \tabularnewline \midrule \midrule
      Dependent Variable: & \multicolumn{6}{c}{Fishing hours}\\
      Model:               & (1)             & (2)             & (3)             & (4)             & (5)             & (6)\\  
      \midrule
      \emph{Variables}\\
      MHW days             & -0.0107$^{***}$ & -0.0066$^{***}$ & -0.0120$^{***}$ & -0.0113$^{***}$ & -0.0113$^{***}$ & -0.0165$^{***}$\\   
                           & (0.0036)        & (0.0024)        & (0.0036)        & (0.0030)        & (0.0036)        & (0.0042)\\   
      SST anomaly          & 0.1224$^{***}$  & 0.1230$^{***}$  & 0.1365$^{***}$  & 0.1308$^{***}$  & 0.1308$^{***}$  & 0.1178$^{***}$\\   
                           & (0.0392)        & (0.0350)        & (0.0393)        & (0.0394)        & (0.0303)        & (0.0386)\\   
      \midrule
      \emph{Fixed-effects}\\
      Cell                 & Yes             & Yes             & Yes             & Yes             & Yes             & Yes\\  
      Month                & Yes             & Yes             & Yes             & Yes             & Yes             & Yes\\  
      Region-Month of year & Yes             & Yes             & Yes             & Yes             & Yes             & Yes\\  
      Region-Year          &                 & Yes             &                 &                 &                 & \\  
      \midrule
      \emph{Fit statistics}\\
      Observations         & 4,533,120       & 686,160         & 486,000         & 686,160         & 686,160         & 686,160\\  
      Squared Correlation  & 0.37703         & 0.52074         & 0.50621         & 0.50691         & 0.50691         & 0.50868\\  
      Pseudo R$^2$         & 0.67998         & 0.74223         & 0.72151         & 0.73454         & 0.73454         & 0.73474\\  
      BIC                  & 22,557,550.6    & 15,074,207.5    & 14,717,132.7    & 15,520,953.8    & 15,520,953.8    & 15,509,086.8\\  
      \midrule \midrule
      \multicolumn{7}{l}{\emph{Signif. Codes: ***: 0.01, **: 0.05, *: 0.1}}\\
   \end{tabular}
\end{table}



\begin{table}[htbp]
\centering
\caption{Sample exclusions: coastal cells and low-coverage cells}
\label{tab:excl}
\begin{tabular}{lrrr}
\hline
model & b & se & p
\\ \hline
excl\_coastal & -0.0115 & 0.0035 & 0.0011 \\
excl\_low\_cov & -0.0115 & 0.0035 & 0.0011 \\
\hline
\end{tabular}
\end{table}


Table~\ref{tab:placebo} and Figure~\ref{fig:perm} report the placebos. The within-region permutation places the observed coefficient outside the 95\% range of \nPermutations{} draws ($p=\pPermutation{}$). The twelve-month lead, however, is significant in the baseline even conditional on lags 0--6 (\bLeadTwelve{}, $p=\pLeadTwelve{}$). This is not anticipation --- the twelve-month autocorrelation of exposure within cell is only 0.05 --- but differential trends: exposure grew more than tenfold over the period, and cells whose exposure grew fastest may also be cells whose effort declined for unrelated reasons (fleet composition, AIS adoption, stock dynamics). Table~\ref{tab:celltrend} adds cell-specific linear trends. The twelve-month lead vanishes (\bLeadTwelveCT{}, $p=\pLeadTwelveCT{}$), the main coefficient falls to \bMhwCellTrend{} ($p=\pMhwCellTrend{}$) and the cumulative six-month effect is no longer significant (\bCumSixCT{}, $p=\pCumSixCT{}$). We therefore report two numbers: the baseline of Table~\ref{tab:main} as the estimate comparable to the literature, and the cell-trend specification as the conservative, placebo-consistent bound. About half of the baseline association is attributable to trends that the fixed effects of equation~\eqref{eq:main} do not remove; the remaining half survives the most demanding specification we can estimate.

\begin{table}[htbp]
\centering
\caption{Placebo tests: 12-month lead and spatial permutation}
\label{tab:placebo}
\begin{tabular}{lrrrr}
\hline
test & term & b & se & p
\\ \hline
placebo\_lead12\_conditional & f(mhw\_days, 12) & -0.0099 & 0.0032 & 0.0020 \\
placebo\_lead12\_unconditional & f(mhw\_days, 12) & -0.0102 & 0.0036 & 0.0047 \\
spatial\_permutation & 50 draws & -0.0113 & 0.0014 & 0.0196 \\
\hline
\end{tabular}
\end{table}


\begin{table}[htbp]
   \caption{\label{tab:celltrend} Cell-specific linear trends: main effect, 12-month lead placebo and lags}
   \centering
   \begin{tabular}{lccc}
      \tabularnewline \midrule \midrule
      Dependent Variable: & \multicolumn{3}{c}{Fishing hours}\\
      Model:               & (1)            & (2)           & (3)\\  
      \midrule
      \emph{Variables}\\
      MHW days             & -0.0046$^{**}$ &               &   \\   
                           & (0.0022)       &               &   \\   
      SST anomaly          & 0.1263$^{***}$ & 0.1164$^{**}$ & 0.1260$^{***}$\\   
                           & (0.0292)       & (0.0480)      & (0.0277)\\   
      f(mhw\_days, 12)     &                & -0.0020       &   \\   
                           &                & (0.0018)      &   \\   
      l(mhw\_days, 0)      &                & -0.0026       & -0.0045$^{***}$\\   
                           &                & (0.0028)      & (0.0017)\\   
      l(mhw\_days, 1)      &                & 0.0029        & 0.0007\\   
                           &                & (0.0027)      & (0.0025)\\   
      l(mhw\_days, 2)      &                & -0.0038$^{*}$ & -0.0039$^{***}$\\   
                           &                & (0.0022)      & (0.0015)\\   
      l(mhw\_days, 3)      &                & 0.0042$^{*}$  & 0.0043$^{*}$\\   
                           &                & (0.0022)      & (0.0023)\\   
      l(mhw\_days, 4)      &                & -0.0032       & -0.0026\\   
                           &                & (0.0029)      & (0.0023)\\   
      l(mhw\_days, 5)      &                & -0.0013       & -0.0004\\   
                           &                & (0.0023)      & (0.0020)\\   
      l(mhw\_days, 6)      &                & 0.0040$^{**}$ & 0.0007\\   
                           &                & (0.0018)      & (0.0021)\\   
      \midrule
      \emph{Fixed-effects}\\
      Cell                 & Yes            & Yes           & Yes\\  
      Month                & Yes            & Yes           & Yes\\  
      Region-Month of year & Yes            & Yes           & Yes\\  
      \midrule
      \emph{Varying Slopes}\\
      time\_id (Cell)      & Yes            & Yes           & Yes\\  
      \midrule
      \emph{Fit statistics}\\
      Observations         & 686,160        & 591,318       & 657,432\\  
      Squared Correlation  & 0.53572        & 0.53244       & 0.53694\\  
      Pseudo R$^2$         & 0.76987        & 0.76121       & 0.77025\\  
      BIC                  & 13,528,193.4   & 11,286,270.8  & 13,318,216.7\\  
      \midrule \midrule
      \multicolumn{4}{l}{\emph{Conley (200km) standard-errors in parentheses}}\\
      \multicolumn{4}{l}{\emph{Signif. Codes: ***: 0.01, **: 0.05, *: 0.1}}\\
   \end{tabular}
   
   \par \raggedright 
   Varying slopes cell x linear time. Conley SE (200 km).
\end{table}




\begin{figure}[htbp]\centering
\includegraphics[width=0.8\linewidth]{figures/fig_placebo_permutation.png}
\caption{Spatial placebo: distribution of the MHW coefficient across \nPermutations{} within-region permutations of the exposure series; the vertical line is the observed coefficient.}\label{fig:perm}
\end{figure}

\section{Information systems contributions}\label{sec:is}
% AUTHOR WRITES. Draft outline of the gains: (1) integration of four previously incompatible sources (GFW, NOAA, IBGE, MHW events) under a single spatial key (H3) and a single temporal key (month), with lineage from raw file to coefficient; (2) reproducibility --- one command rebuilds every table, figure and number, package versions are pinned, downloads carry checksums; (3) resilience --- block-level checkpoints make multi-hour computations restartable; (4) extensibility --- productivity, protected areas and ports enter as additional layers without changing the analytical core; (5) monitoring --- the same cube feeds a dashboard for period/region selection, anomaly detection and export of aggregates permitted by the data licences.

\section{Economic and policy implications}\label{sec:policy}
% AUTHOR WRITES. Draft outline: the threshold nature of the response (20+ days) implies that adaptation instruments should be triggered by duration, not by anomaly; the absence of nearby refuge and the spread to 60--150 km imply higher steaming costs and changes in the spatial footprint that closed-season and MPA design should anticipate; the regional gradient (North most exposed) identifies where compensatory instruments such as Seguro-Defeso would be most relevant; the AIS-tracked fleet is the industrial segment, so distributional conclusions for artisanal communities require other data.

\section{Limitations}\label{sec:limits}
The results describe the AIS-tracked fleet, dominated by drifting longliners; they say nothing about the artisanal fleet, which is invisible to AIS. Apparent fishing hours are a classifier output with fleet-specific error \citep{hintzen2025}. SST at $0.25^{\circ}$ is coarse near the coast, and the H3 cell is larger than the SST pixel, so exposure is measured with error that attenuates the estimates. The design is associational: the placebos exclude anticipation and, under cell trends, differential trends, but time-varying unobserved shocks correlated with exposure within cell remain possible. The analysis is monthly; sub-monthly responses are not observed. Ocean productivity, protected areas and port distance were not included and are left for future work.

\section{Conclusion}\label{sec:conclusion}
% AUTHOR WRITES.

% data_code_availability.tex — repository infrastructure statement.
% Fill every [bracket] before submission and re-check each licence page.

\section*{Data and code availability}
All code that collects, processes and analyses the data, together with the
reproducibility documentation, is publicly available at
\url{https://github.com/valbergregory/marine-fisheries-data-cube}
(release [tag], commit [hash]). The analysis is orchestrated by a
\texttt{targets} pipeline; package versions are pinned in \texttt{renv.lock};
every downloaded file is registered with a checksum in
\texttt{data/metadata/checksums.csv}. Every table, figure and number reported
in this article is produced by that pipeline; none was typed by hand.

Raw data are \emph{not} redistributed in the repository and remain governed by
their providers' terms. Apparent fishing effort comes from Global Fishing Watch
(4Wings API v3, accessed [date]; data licensed CC BY-SA 4.0, free registration
required). Sea-surface temperature comes from the NOAA OISST v2.1 daily product
(NOAA PSL, accessed [date]; United States government data, public domain).
Territorial layers come from IBGE (public data). The full inventory, including
licences, access dates and known limitations, is in
\texttt{config/data\_sources.yml} and \texttt{docs/data\_inventory.md}.

Derived aggregates released with this article are shared under
[CC BY-SA 4.0, to match the Global Fishing Watch share-alike requirement],
after verification of the provider's API terms of service.

% ai_disclosure.tex — paste into main.tex immediately before \bibliography (Elsevier
% placement) or inside the Methods/Acknowledgements section (Springer, Wiley, OUP, CUP,
% T&F, SAGE). Fill every [bracket]. Keep only ONE language block.
% The same text goes into the submission form and the cover letter.

% ---------- English ----------
\section*{Declaration of generative AI and AI-assisted technologies}
During the preparation of this work the author used Claude Code (Anthropic; model
Claude Fable 5.1; desktop application version [X.Y.Z]) to assist in writing, testing
and documenting the data-collection, data-processing and estimation code released in
the public repository, and to prepare the reproducibility documentation of that
repository. The author designed the study, specified every procedure, reviewed and
executed all code, and verified all outputs.
% Keep the next sentence only if an LLM is part of the METHOD (articles 3, 4, 5):
An open-weights language model ([name, version, hash]) was used as a measurement
instrument, as described in Section~\ref{sec:methods}; its outputs were validated
against human annotation and did not replace the author's analysis.
No generative AI tool was used to write, draft or paraphrase the text of this article.
% Keep the next sentence only if a language-editing tool was used on author-written text:
A language-editing tool ([name, version]) was used to check grammar and style of
author-written text.
After using these tools the author reviewed and edited the content as needed and takes
full responsibility for the content of the article.

% ---------- Português ----------
% \section*{Declaração sobre o uso de inteligência artificial generativa}
% Na preparação deste trabalho o autor utilizou o Claude Code (Anthropic; modelo Claude
% Fable 5.1; aplicativo versão [X.Y.Z]) como apoio à escrita, ao teste e à documentação
% do código de coleta, tratamento e estimação disponibilizado no repositório público, e
% à documentação de reprodutibilidade desse repositório. A concepção do estudo, a
% especificação de cada procedimento, a revisão e a execução de todo o código e a
% verificação de todos os resultados foram realizadas pelo autor.
% % Manter só se um LLM integra o MÉTODO (artigos 3, 4, 5):
% Um modelo de linguagem de pesos abertos ([nome, versão, hash]) foi empregado como
% instrumento de medida, conforme descrito na Seção~\ref{sec:methods}; seus resultados
% foram validados contra anotação humana e não substituíram a análise do autor.
% Nenhuma ferramenta de IA generativa foi usada para redigir, esboçar ou parafrasear o
% texto deste artigo.
% % Manter só se houve revisão linguística automática de texto escrito pelo autor:
% Uma ferramenta de revisão linguística ([nome, versão]) foi usada para verificar
% gramática e estilo de texto escrito pelo autor.
% Após o uso dessas ferramentas o autor revisou e editou o conteúdo conforme necessário
% e assume integral responsabilidade pelo conteúdo do artigo (Portaria CNPq
% 2.664/2026, art.~9\textsuperscript{o}, I, ``c'' e ``d'').


\bibliographystyle{apalike}
\bibliography{references}
\end{document}
